% ============================================================
%  Preámbulo: paquetes, estilo y macros del informe
% ============================================================

% --- Idioma y codificación ---
\usepackage[utf8]{inputenc}
\usepackage[T1]{fontenc}
\usepackage[spanish,es-noshorthands,es-tabla]{babel}
\usepackage{lmodern}
\usepackage{microtype}

% --- Página ---
\usepackage[margin=2.5cm]{geometry}
\usepackage{fancyhdr}
\usepackage{titlesec}
\setlength{\parskip}{0.5em}
\setlength{\parindent}{0pt}

% --- Matemática ---
\usepackage{amsmath,amssymb,amsthm,mathtools}
\usepackage{bm}

% --- Algoritmos (pseudocódigo) ---
\usepackage[ruled,vlined,linesnumbered,spanish,onelanguage]{algorithm2e}

% --- Código Python ---
\usepackage{listings}
\usepackage{xcolor}

% --- Figuras, tablas y árboles de recursión ---
\usepackage{graphicx}
\usepackage{booktabs}
\usepackage{array}
\usepackage{float}
\usepackage{enumitem}
\usepackage{caption}
\usepackage{subcaption}
\usepackage{tikz}
\usetikzlibrary{trees,positioning,arrows.meta,calc}
\usepackage{forest}

% --- Cajas para resultados ---
\usepackage[most]{tcolorbox}

% --- Enlaces (cargar al final) ---
\usepackage[hidelinks]{hyperref}
\usepackage[nameinlink,spanish]{cleveref}

\graphicspath{{figuras/}}

% ============================================================
%  Encabezado y pie
% ============================================================
\pagestyle{fancy}
\fancyhf{}
\fancyhead[L]{\small\curso}
\fancyhead[R]{\small\tituloinforme}
\fancyfoot[C]{\thepage}
\renewcommand{\headrulewidth}{0.4pt}

\titleformat{\section}{\Large\bfseries}{\thesection}{0.8em}{}
\titleformat{\subsection}{\large\bfseries}{\thesubsection}{0.8em}{}

% ============================================================
%  Estilo para código Python
% ============================================================
\definecolor{codefondo}{RGB}{247,247,247}
\definecolor{codeclave}{RGB}{0,80,160}
\definecolor{codecoment}{RGB}{110,110,110}
\definecolor{codetexto}{RGB}{160,40,40}
\definecolor{codenum}{RGB}{150,150,150}

\lstdefinestyle{python}{
  language=Python,
  backgroundcolor=\color{codefondo},
  basicstyle=\ttfamily\small,
  keywordstyle=\color{codeclave}\bfseries,
  commentstyle=\color{codecoment}\itshape,
  stringstyle=\color{codetexto},
  numbers=left,
  numberstyle=\tiny\color{codenum},
  numbersep=8pt,
  frame=single,
  rulecolor=\color{black!15},
  breaklines=true,
  showstringspaces=false,
  tabsize=4,
  columns=fullflexible,
  keepspaces=true,
  xleftmargin=1.5em,
  framexleftmargin=1.5em,
  % Acentos y ñ dentro del código
  literate=
    {á}{{\'a}}1 {é}{{\'e}}1 {í}{{\'i}}1 {ó}{{\'o}}1 {ú}{{\'u}}1
    {Á}{{\'A}}1 {É}{{\'E}}1 {Í}{{\'I}}1 {Ó}{{\'O}}1 {Ú}{{\'U}}1
    {ñ}{{\~n}}1 {Ñ}{{\~N}}1 {ü}{{\"u}}1 {¿}{{?`}}1 {¡}{{!`}}1
}
\lstset{style=python}
\renewcommand{\lstlistingname}{Código}

% ============================================================
%  Macros de notación asintótica
% ============================================================
\newcommand{\T}[1]{\Theta\!\left(#1\right)}
\newcommand{\Oh}[1]{O\!\left(#1\right)}
\newcommand{\Om}[1]{\Omega\!\left(#1\right)}
\newcommand{\oh}[1]{o\!\left(#1\right)}
\newcommand{\om}[1]{\omega\!\left(#1\right)}
\newcommand{\N}{\mathbb{N}}
\newcommand{\R}{\mathbb{R}}
\DeclarePairedDelimiter{\ceil}{\lceil}{\rceil}
\DeclarePairedDelimiter{\floor}{\lfloor}{\rfloor}
\DeclarePairedDelimiter{\norma}{\lVert}{\rVert}

% ============================================================
%  Entornos teóricos
% ============================================================
\theoremstyle{definition}
\newtheorem{definicion}{Definición}[section]
\newtheorem{proposicion}{Proposición}[section]
\newtheorem{lema}{Lema}[section]
\theoremstyle{remark}
\newtheorem*{observacion}{Observación}
\renewcommand{\qedsymbol}{$\blacksquare$}

% ============================================================
%  Estructura de cada ejercicio
% ============================================================
% Enunciado del ejercicio
\newtcolorbox{enunciado}{
  colback=black!3, colframe=black!35, boxrule=0.5pt, arc=2pt,
  left=6pt, right=6pt, top=4pt, bottom=4pt,
  fonttitle=\bfseries, title=Enunciado
}
\newcommand{\paso}[2]{\par\medskip\noindent\textbf{Paso #1. #2.}\par\nopagebreak[4]}

% Resultado final destacado
\newtcolorbox{resultado}{
  colback=blue!4, colframe=blue!50!black, boxrule=0.6pt, arc=2pt,
  left=6pt, right=6pt, top=4pt, bottom=4pt,
  fonttitle=\bfseries, title=Resultado
}

% Intuición / comentario pedagógico
\newtcolorbox{intuicion}{
  colback=green!4, colframe=green!40!black, boxrule=0.5pt, arc=2pt,
  left=6pt, right=6pt, top=4pt, bottom=4pt,
  fonttitle=\bfseries, title=Intuición
}
